% ============================================================================
% CHEATSHEET TEMPLATE
% ============================================================================
% Based on the ANA-G cheatsheet style
%
% Features:
% - 4-column landscape layout on A4
% - Colored boxes for main concepts and sub-concepts
% - Compact formatting optimized for information density
% - Mathematical typesetting support
%
% Usage:
% 1. Replace metadata (title, author, date)
% 2. Adjust number of columns if needed (see \Ncolumn)
% 3. Add your content using mainbox and subbox environments
% 4. Compile with pdflatex (run twice for proper layout)
% ============================================================================

\documentclass[a4paper, fontsize=7pt]{scrartcl}

% ============================================================================
% ENCODING & LANGUAGE
% ============================================================================
\usepackage[utf8]{inputenc}
% \usepackage[german]{babel}    % Uncomment for German hyphenation
% \usepackage[english]{babel}   % Uncomment for English hyphenation

% ============================================================================
% PAGE LAYOUT
% ============================================================================
% 4-column landscape layout with minimal margins
\usepackage[landscape, left=0.75cm, top=1cm, right=0.75cm, bottom=1.5cm, footskip=5pt]{geometry}
\usepackage{flowfram}
\ffvadjustfalse
\setlength{\columnsep}{1cm}

% Number of columns (adjust as needed: 3-4 recommended for A4)
\Ncolumn{4}

% Remove paragraph indentation
\setlength{\parindent}{0cm}

% ============================================================================
% COLORS & BOXES
% ============================================================================
\usepackage[most]{tcolorbox}

% Base style for all boxes
\tcbset{
  base/.style={
    boxrule=0.3mm,
    leftrule=1mm,
    left=1.75mm,
    arc=0mm,
    fonttitle=\bfseries,
    colbacktitle=black!10!white,
    coltitle=black,
    toptitle=0.75mm,
    bottomtitle=0.25mm,
    title={#1}
  }
}

% Main concept box (green accent)
% Use for: definitions, theorems, important formulas
\definecolor{maincolor}{rgb}{0.004, 0.5, 0.35}  % Nice green
\newtcolorbox{mainbox}[1]{
  colframe=maincolor,
  base={#1}
}

% Sub concept box (gray accent)
% Use for: examples, remarks, minor concepts
\newtcolorbox{subbox}[1]{
  colframe=black!20!white,
  base={#1}
}

% Alternative colors (uncomment to use)
% \definecolor{maincolor}{rgb}{0.2, 0.4, 0.8}   % Blue
% \definecolor{maincolor}{rgb}{0.8, 0.2, 0.2}   % Red
% \definecolor{maincolor}{rgb}{0.5, 0.2, 0.6}   % Purple

% ============================================================================
% MATHEMATICS
% ============================================================================
\usepackage{amsmath, amsthm, mathtools, amssymb}
\usepackage{amsfonts, amstext}
\allowdisplaybreaks  % Allow page breaks in equations

% Common math shortcuts (customize as needed)
\def\limn{\lim_{n\to \infty}}
\def\limxo{\lim_{x\to 0}}
\def\limxi{\lim_{x\to\infty}}
\def\limxn{\lim_{x\to-\infty}}
\def\sumk{\sum_{k=1}^\infty}
\def\sumn{\sum_{n=0}^\infty}

% Number sets
\def\N{\mathbb{N}}  % Natural numbers
\def\Z{\mathbb{Z}}  % Integers
\def\Q{\mathbb{Q}}  % Rationals
\def\R{\mathbb{R}}  % Reals
\def\C{\mathbb{C}}  % Complex

% Calculus
\def\dx{\,dx}
\def\dt{\,dt}
\def\dy{\,dy}
\newcommand{\ddx}[1]{\frac{d#1}{dx}}
\newcommand{\ddn}[2]{\frac{d#1}{d#2}}
\newcommand{\pdn}[2]{\frac{\partial #1}{\partial #2}}

% Big O notation
\def\oof{\mathcal{O}}

% ============================================================================
% TABLES
% ============================================================================
\usepackage{tabularx, multirow, booktabs}
\renewcommand*{\arraystretch}{1.5}  % More spacing in tables

% ============================================================================
% LISTS
% ============================================================================
\usepackage{enumitem}
\setitemize{itemsep=0.5pt, topsep=2pt}
\setenumerate{itemsep=0.75pt, topsep=2pt}

% ============================================================================
% GRAPHICS & TIKZ
% ============================================================================
\usepackage{graphicx}
\usepackage{tikz}
\usetikzlibrary{
  positioning,
  calc,
  fit,
  matrix,
  decorations,
  angles,
  quotes,
  babel,
  3d,
  arrows.meta,
  shadings,
  shapes.geometric
}

\tikzset{
  >=stealth,
  % Add custom TikZ styles here
}

% ============================================================================
% HYPERLINKS
% ============================================================================
\usepackage{hyperref}
\hypersetup{
  colorlinks=true,
  linkcolor=maincolor,
  urlcolor=maincolor
}

% ============================================================================
% ADDITIONAL PACKAGES
% ============================================================================
\usepackage{centernot}   % For \centernot command
\usepackage{marvosym}    % Additional symbols
\usepackage{wasysym}     % Additional symbols

% ============================================================================
% DOCUMENT METADATA
% ============================================================================
\title{Cheatsheet DMATH-CODE}
\author{Thomas Schnyder}
\date{September 2026}

% Table of contents depth
\setcounter{tocdepth}{4}
\setcounter{secnumdepth}{4}

% ============================================================================
% DOCUMENT START
% ============================================================================
\begin{document}

\maketitle

% ============================================================================
% SECTION 1: WAHRSCHEINLICHKEIT
% ============================================================================
\section{Wahrscheinlichkeit}

\subsection{Zufallsexperiment \& Ergebnismenge}

\begin{mainbox}{Zufallsexperiment}
Ein Vorgang, der
\begin{itemize}
    \item beliebig oft unter gleichen Bedingungen wiederholt werden kann und
    \item dessen Ergebnis nicht mit Sicherheit vorhergesagt werden kann.
\end{itemize}
\textbf{Ergebnismenge} $\Omega$: Menge \emph{aller} möglichen Ergebnisse
(bekannt!).\\
Mehrere Teilexperimente $\Rightarrow$ kartesisches Produkt, z.B.
$\Omega = \{0,1\}^n$, $|\Omega| = 2^n$.
\end{mainbox}

\begin{subbox}{Beispiele Ergebnismenge}
\begin{itemize}
    \item 1 Würfel: $\Omega = \{1,\ldots,6\}$
    \item 2 Würfel: $\Omega = \{1,\ldots,6\}^2 = \{11, 12, \ldots, 66\}$,
          $|\Omega| = 36$
    \item 3 Münzen (CHF, \$, £): $\{K,Z\}^3$, $|\Omega| = 8$
    \item Tennis (Sieger: zuerst 2 in Folge oder 3 total): $\Omega = \{$AA,
          BB, ABB, BAA, ABAA, BABB, ABABA, ABABB, BABAB, BABAA$\}$,
          $|\Omega| = 10$
\end{itemize}
\end{subbox}

\subsection{Wahrscheinlichkeitsfunktion}

\begin{mainbox}{Wahrscheinlichkeitsfunktion $P$}
$P$ ordnet jedem Ergebnis $\omega \in \Omega = \{\omega_1, \ldots, \omega_n\}$
eine reelle Zahl $P(\omega) \in [0,1]$ zu, sodass
$$\boxed{\sum_{\omega \in \Omega} P(\omega) = P(\omega_1) + \ldots + P(\omega_n) = 1}$$
\textbf{Uniform:} alle Ergebnisse gleich wahrscheinlich:
$$P(\omega_1) = \ldots = P(\omega_n) = \frac{1}{|\Omega|}$$
\end{mainbox}

\begin{subbox}{Beispiele}
\begin{itemize}
    \item Fairer Würfel: $P(k) = \frac{1}{6}$; 2 Würfel: $P(\omega) = \frac{1}{36}$
    \item Münze mit $P(K) = 0.3$ $\Rightarrow$ $P(Z) = 1 - 0.3 = 0.7$
    \item Würfel, jede Augenzahl 10\,\% wahrscheinlicher als die vorherige:
          $P(1) = x$, $\sum_{k=0}^{5} 1.1^k x = 7.71561\,x = 1$
          $\Rightarrow P(1) \approx 0.1296$, \ldots, $P(6) \approx 0.2087$
\end{itemize}
\end{subbox}

\begin{subbox}{Beispiel: Zufallsgenerator 0--7 aus Bit-Generator}
\textbf{Var. 1:} 3 Bits als Binärzahl: $\Omega = \{0,1\}^3$, uniform
$P(k) = \frac{1}{8}$.\\
\textbf{Var. 2:} Anzahl Nullen unter 7 Bits: $|\Omega| = 2^7$,
$$P(k) = \binom{7}{k} \cdot \frac{1}{2^7} \quad \text{(nicht uniform!)}$$
$P(0) = P(7) \approx 0.0078$, $P(1) = P(6) \approx 0.0547$,\\
$P(2) = P(5) \approx 0.1641$, $P(3) = P(4) \approx 0.2734$.
\end{subbox}

\subsection{Ereignisse}

\begin{mainbox}{Ereignis}
Jede Teilmenge $E = \{\omega_1, \ldots, \omega_k\} \subseteq \Omega$ heisst
\textbf{Ereignis}.
$$\boxed{P(E) = \sum_{\omega \in E} P(\omega)}$$
\textbf{Uniform (Laplace):} \ $P(E) = \dfrac{|E|}{|\Omega|}$
\ (günstige / mögliche)
\end{mainbox}

\begin{mainbox}{Rechenregeln (Satz 1)}
\begin{itemize}
    \item Sicheres Ereignis: $P(\Omega) = 1$ (trifft immer zu)
    \item Unmögliches Ereignis: $P(\{\}) = 0$ (trifft nie zu)
    \item Gegenereignis: $P(\overline{A}) = 1 - P(A)$
    \item Vereinigung:
          $P(A \cup B) = P(A) + P(B) - P(A \cap B)$
\end{itemize}
\end{mainbox}

\begin{subbox}{Beispiel Roulette (38 Felder)}
$P(\overline{\text{odd}}) = 1 - \frac{18}{38} = \frac{20}{38}$ \ (0 und 00!)\\
$P(\text{odd} \cup \text{red}) = \frac{18}{38} + \frac{18}{38} - \frac{10}{38}
= \frac{26}{38}$
\end{subbox}

\begin{subbox}{Denkprozess: ``mindestens ein\ldots''}
\begin{enumerate}
    \item Direkt zählen ist mühsam (viele Fälle).
    \item Gegenereignis ``kein \ldots'' ist \emph{ein} einfacher Fall.
    \item $P(\text{mind. ein}) = 1 - P(\text{kein})$
\end{enumerate}
\textbf{8-Bit-Wort mit mind. einer 1:} $1 - \frac{1}{2^8} \approx 0.9961$
\end{subbox}

\subsection{Typische Aufgaben}

\begin{subbox}{Chevalier de Méré}
\textbf{Spiel 1:} 4 Würfe, mind. eine Sechs:
$$1 - \left(\tfrac{5}{6}\right)^4 \approx 0.5177$$
\textbf{Spiel 2:} 24 Würfe mit 2 Würfeln, mind. eine Doppelsechs:
$$1 - \left(\tfrac{35}{36}\right)^{24} \approx 0.4914$$
$\Rightarrow$ Die Chancen sind \emph{nicht} gleich, obwohl $\frac{1}{6} \cdot 4 =
\frac{1}{36} \cdot 24$.
\end{subbox}

\begin{subbox}{Teilbarkeit (Inklusion--Exklusion)}
Zahl uniform aus $1..100$, teilbar durch 2 oder 5:
$$\tfrac{50}{100} + \tfrac{20}{100} - \tfrac{10}{100} = 0.6$$
(Schnitt = durch 10 teilbar)
\end{subbox}

\begin{subbox}{Kein Index-Treffer}
5 Zufallszahlen aus $\{0,\ldots,4\}$ in Liste; kein Element $=$ Index:
jede Stelle hat 4 erlaubte Werte $\Rightarrow$
$$P = \frac{4^5}{5^5} = 0.32768$$
\end{subbox}

\begin{subbox}{Gruppen gleich gemischt}
6 Frauen, 6 Männer in 2 gleich grosse Gruppen:
$|\Omega| = \binom{12}{6} = 924$
$$P = \binom{6}{3}\binom{6}{3} \cdot \frac{1}{924} \approx 0.4329$$
\end{subbox}

\begin{subbox}{Geburtstagsproblem-Typ}
5 Personen merken sich je eine Ziffer 0--9; mind. 2 gleich:
$$1 - \frac{10 \cdot 9 \cdot 8 \cdot 7 \cdot 6}{10^5} = 0.6976$$
15 Jobs auf 20 Server, alle auf verschiedenen Servern:
$$\frac{20!}{5!} \cdot \frac{1}{20^{15}} \approx 0.00062$$
\end{subbox}

\begin{subbox}{Linear Probing (Array Länge 6)}
2 Elemente uniform; Kollision $\to$ nächste freie Stelle (zyklisch).
Benachbart, falls 2.\ Element auf $i-1$, $i+1$ \emph{oder} $i$ (dann $\to i+1$)
fällt:
$$P = \frac{3}{6} = \frac{1}{2}$$
\end{subbox}

\begin{subbox}{Häufige Fehler \& Tipps}
\begin{itemize}
    \item ``mindestens eins'' $\Rightarrow$ fast immer Gegenereignis.
    \item Bei Vereinigung den Schnitt abziehen (Doppelzählung).
    \item Nicht jedes Experiment ist uniform (z.B. Anzahl Nullen, Summe von
          Würfeln) $\Rightarrow$ Laplace nur, wenn alle $\omega$ gleich
          wahrscheinlich.
    \item Wahrscheinlichkeit $\cdot$ Anzahl Versuche ist \emph{keine}
          Wahrscheinlichkeit (de Méré).
    \item Ordnung in $\Omega$ beachten: $(1,2) \neq (2,1)$ bei 2 Würfeln.
\end{itemize}
\end{subbox}

% ============================================================================
% SECTION 2: BEDINGTE WAHRSCHEINLICHKEIT
% ============================================================================
\section{Bedingte Wahrscheinlichkeit}

\subsection{Definition \& Produktregel}

\begin{mainbox}{Bedingte Wahrscheinlichkeit}
Ist bekannt, dass nur noch Ergebnisse aus $E \subseteq \Omega$ vorkommen
($P(E) > 0$):
$$\boxed{P(A \mid E) = \frac{P(A \cap E)}{P(E)}}$$
\begin{itemize}
    \item $P(\cdot \mid E)$ gibt nur $\omega \in E$ positive W'keit;
          $\sum_{\omega \in E} P(\omega \mid E) = 1$.
    \item \textbf{Produktregel:} $P(A \cap E) = P(E) \cdot P(A \mid E)$
    \item Uniform: $P(A \mid E) = \dfrac{|A \cap E|}{|E|}$ \ ($E$ wird
          neues $\Omega$)
\end{itemize}
\end{mainbox}

\begin{subbox}{Beispiel Zufallszahl 1--10, $E = $ gerade}
$P(2) = P(E) \cdot P(2 \mid E) = \frac{5}{10} \cdot \frac{1}{5}$
$\Rightarrow P(2 \mid E) = \frac{1}{5}$\\
$P(\text{3er-Zahl}) = \frac{3}{10}$, \ $P(\text{3er-Zahl} \mid E) =
\frac{1}{5}$ \ (verändert)\\
$P(\text{5er-Zahl}) = \frac{2}{10}$, \ $P(\text{5er-Zahl} \mid E) =
\frac{1}{5}$ \ (\emph{gleich!})
\end{subbox}

\begin{mainbox}{Unabhängigkeit}
$A$ und $E$ heissen \textbf{unabhängig}, wenn
$$P(A \mid E) = P(A)$$
(Information $E$ ändert nichts an $A$). Äquivalent:
$$\boxed{P(A \cap E) = P(A) \cdot P(E)}$$
Symmetrisch: $P(E \mid A) = \frac{P(A \cap E)}{P(A)} =
\frac{P(A)P(E)}{P(A)} = P(E)$.
\end{mainbox}

\begin{subbox}{Beispiel 3 Würfel: Wette auf $A$ oder $B$?}
$A$: Summe $\le 9$; \ $B$: mind.\ 2 gleiche Augenzahlen.\\
$P(A) = \frac{81}{216} = 0.375$, \
$P(B) = 1 - \frac{6 \cdot 5 \cdot 4}{216} = \frac{96}{216} \approx 0.444$
$\Rightarrow$ auf $B$ setzen.\\
\textbf{Mit Info} $E$: mind.\ eine Eins, $|E| = 216 - 5^3 = 91$:\\
$|A \cap E| = 81 - 20 = 61$ (20: Summe $\le 9$ ohne Eins),\\
$|B \cap E| = 91 - 3 \cdot 5 \cdot 4 = 31$\\
$P(A \mid E) = \frac{61}{91} \approx 0.670$, \
$P(B \mid E) = \frac{31}{91} \approx 0.341$ $\Rightarrow$ \textbf{auf $A$
wechseln!}
\end{subbox}

\begin{subbox}{Beispiel Unabhängigkeit: roter \& blauer Würfel}
$A$: rot ungerade, $B$: Summe ungerade.
$P(A) = P(B) = \frac{1}{2}$,\\
$A \cap B$ = rot ungerade \& blau gerade: $\frac{9}{36} = \frac{1}{4} =
P(A)P(B)$ $\Rightarrow$ \textbf{unabhängig}.
\end{subbox}

\begin{subbox}{Beispiel 10 Münzwürfe, $H$: mind.\ 9 Mal Kopf}
$P(H) = \frac{\binom{10}{9} + \binom{10}{10}}{2^{10}} = \frac{11}{1024}$
\begin{itemize}
    \item 1.\ Wurf Kopf: Rest 9 mit $\ge 8$ K: $\frac{10}{2^9} \approx 0.0195$
    \item 1.\ Wurf Zahl: Rest 9 alle K: $\frac{1}{2^9} \approx 0.00195$
    \item ersten 3 Kopf: Rest 7 mit $\ge 6$ K: $\frac{8}{2^7} = \frac{1}{16}$
    \item mind.\ 8 Köpfe: $\frac{11}{45 + 10 + 1} = \frac{11}{56} \approx
          0.196$
\end{itemize}
\end{subbox}

\subsection{Mehrstufige Zufallsexperimente}

\begin{mainbox}{Wahrscheinlichkeitsbaum}
Zufallsexperiment aus mehreren gekoppelten Stufen.
\begin{itemize}
    \item \textbf{Knoten}: Zwischenstände; \textbf{Blätter}: Ergebnisse
    \item \textbf{Kanten}: bedingte W'keiten $P(\text{Knoten } B \mid
          \text{Knoten } A)$; Kanten aus einem Knoten summieren zu 1
    \item \textbf{Pfadregel}: $P(\text{Blatt})$ = Produkt der Kanten
          entlang des Pfads (Produktregel)
    \item \textbf{Summenregel}: $P(\text{Ereignis})$ = Summe der
          zugehörigen Pfade
\end{itemize}
\end{mainbox}

\begin{subbox}{Beispiel Tennis, $P(\text{A gewinnt Spiel}) = p = 0.6$}
\begin{center}
\begin{tikzpicture}[level distance=7mm,
  level 1/.style={sibling distance=18mm},
  level 2/.style={sibling distance=9mm},
  every node/.style={font=\scriptsize, inner sep=1pt},
  edge from parent/.style={draw, ->}]
\node {start}
  child {node {A}
    child {node {AA} edge from parent node[left] {$p$}}
    child {node {AB\ldots} edge from parent node[right] {$q$}}
    edge from parent node[left] {$p$}}
  child {node {B}
    child {node {BA\ldots} edge from parent node[left] {$p$}}
    child {node {BB} edge from parent node[right] {$q$}}
    edge from parent node[right] {$q$}};
\end{tikzpicture}
\end{center}
A gewinnt: AA, BAA, ABAA, ABABA, BABAA ($q = 0.4$):
\begin{align*}
P &= p^2 + qp^2 + pqp^2 + 2p^3q^2\\
  &= 0.36 + 0.144 + 0.0864 + 0.06912\\
  &\approx 0.6595
\end{align*}
\end{subbox}

\begin{subbox}{Beispiel Wetterregel (So.\ trocken)}
T$\to$T $\frac{5}{6}$, T$\to$N $\frac{1}{6}$; \ N$\to$N $\frac{2}{3}$,
N$\to$T $\frac{1}{3}$. Stufenweise (Summenregel):\\
Mo: $T = \frac{5}{6}$, $N = \frac{1}{6}$\\
Di: $T = \frac{5}{6}\cdot\frac{5}{6} + \frac{1}{6}\cdot\frac{1}{3} =
\frac{3}{4}$\\
Mi: $T = \frac{3}{4}\cdot\frac{5}{6} + \frac{1}{4}\cdot\frac{1}{3} =
\frac{17}{24} \approx 0.708$\\
Mind.\ einmal nass bis Mi: $1 - \left(\frac{5}{6}\right)^3 =
\frac{91}{216} \approx 0.421$
\end{subbox}

\begin{subbox}{Weitere typische Aufgaben}
\textbf{Geburtstagsproblem:} kleinstes $n$ mit $P(\text{2 gleich}) \ge 50\%$:
$$1 - \prod_{i=0}^{n-1} \frac{365 - i}{365} \ge 0.5 \Rightarrow
\boxed{n = 23}\ (0.507)$$
\textbf{$n$ Dateien auf $n$ Server}, jeder passend:
$P = \frac{n!}{n^n}$\\
\textbf{Best-of-five}, $p = 0.45$ (3 Siege vor 3 Niederlagen):
$$P = p^3\left(1 + 3q + 6q^2\right) \approx 0.407$$
($\binom{2+k}{k}$ Reihenfolgen mit $k$ Niederlagen vor dem 3.\ Sieg)
\end{subbox}

\begin{subbox}{Häufige Fehler \& Tipps}
\begin{itemize}
    \item $P(A \mid E) \neq P(E \mid A)$ -- Reihenfolge beachten!
    \item Unabhängig $\neq$ disjunkt: disjunkte $A, E$ mit $P > 0$ sind
          \emph{abhängig} ($P(A \cap E) = 0$).
    \item Unabhängigkeit immer \emph{nachrechnen}
          ($P(A\cap E) = P(A)P(E)$), nicht raten.
    \item Im Baum: Kanten sind \emph{bedingte} W'keiten; multiplizieren
          entlang Pfad, addieren über Pfade.
    \item Bei Bedingung $E$: einfach $E$ als neues $\Omega$ behandeln.
\end{itemize}
\end{subbox}

% ============================================================================
% SECTION 3: SATZ VON BAYES
% ============================================================================
\section{Satz von Bayes}

\subsection{Satz von Bayes \& totale Wahrscheinlichkeit}

\begin{mainbox}{Satz von Bayes}
Für beliebige Ereignisse $A, B$ mit $P(B) > 0$:
$$\boxed{P(A \mid B) = \frac{P(A) \cdot P(B \mid A)}{P(B)}}$$
Dreht die Bedingung um: aus $P(B \mid A)$ (oft bekannt, z.B.\
Testgüte) wird $P(A \mid B)$ (gesucht, z.B.\ ``krank, wenn positiv'').\\
Herleitung: Definition + Produktregel im Zähler,
$P(A \cap B) = P(A) \cdot P(B \mid A)$.
\end{mainbox}

\begin{mainbox}{Totale Wahrscheinlichkeit (Partition)}
Ist $\Omega = \Omega_1 \cup \ldots \cup \Omega_n$ eine \textbf{Partition}
(paarweise disjunkt, decken $\Omega$ ab):
$$\boxed{P(B) = \sum_{i=1}^{n} P(\Omega_i) \cdot P(B \mid \Omega_i)}$$
Meist $\Omega_1 = A$, $\Omega_2 = \overline{A}$:
$$P(A|B) = \frac{P(A)P(B|A)}{P(A)P(B|A) + P(\bar{A})P(B|\bar{A})}$$
Im Baum: Zähler = \emph{ein} Pfad, Nenner = \emph{alle} Pfade zu $B$.
\end{mainbox}

\begin{subbox}{Beispiel Zufallszahl 1--10, $G$ = gerade, $U$ = ungerade}
$P(\text{3er}) = P(G)P(\text{3er} \mid G) + P(U)P(\text{3er} \mid U)$\\
$\phantom{P(\text{3er})} = \frac{5}{10}\cdot\frac{1}{5} +
\frac{5}{10}\cdot\frac{2}{5} = \frac{3}{10}$\\
$P(G \mid \text{3er}) = \dfrac{\frac{5}{10} \cdot \frac{1}{5}}{\frac{3}{10}}
= \frac{1}{3}$
\end{subbox}

\begin{mainbox}{Subjektive Interpretation}
\textbf{a-priori}-Überzeugung $P(A)$ wird durch Beobachtung $B$ zur
\textbf{a-posteriori}-Überzeugung $P(A \mid B)$ aktualisiert.
Iterativ: Posterior wird neuer Prior für nächste Beobachtung.\\
Anwendungen: Bayes-Klassifikatoren (Spam), Diagnostik, Fehlersuche,
Spracherkennung, Robotik.
\end{mainbox}

\begin{subbox}{Denkprozess: Bayes-Aufgabe}
\begin{enumerate}
    \item Ereignisse benennen: Ursache $A$ (verborgen), Beobachtung $B$.
    \item Gegeben sortieren: Prior $P(A)$, Likelihoods $P(B \mid A)$,
          $P(B \mid \overline{A})$.
    \item Baum zeichnen: 1.\ Stufe Ursache, 2.\ Stufe Beobachtung.
    \item $P(B)$ via Partition (alle Pfade zu $B$).
    \item Bayes: $P(A \mid B) = \frac{\text{Pfad } A \to B}{P(B)}$.
\end{enumerate}
\end{subbox}

\subsection{Typische Aufgaben}

\begin{subbox}{Ziegenproblem (Let's Make a Deal)}
Wahl Tor 1, Moderator öffnet Tor 3. $G_i$: Gewinn hinter $i$,
$M_j$: öffnet $j$. $P(M_3 \mid G_1) = \frac{1}{2}$,
$P(M_3 \mid G_2) = 1$, $P(M_3 \mid G_3) = 0$:
$$P(G_2 \mid M_3) = \frac{\frac{1}{3} \cdot 1}{\frac{1}{3}\cdot\frac{1}{2}
+ \frac{1}{3} \cdot 1 + \frac{1}{3} \cdot 0} = \frac{2}{3}$$
$\Rightarrow$ \textbf{Wechseln} verdoppelt die Gewinnchance.
\end{subbox}

\begin{subbox}{Urnen: woher stammt die Kugel?}
Urne $A$: Nr.\ 1--9, Urne $B$: Nr.\ 1--5; Urne uniform, gerade Nr.\ $G$:
$$P(A \mid G) = \frac{\frac{1}{2}\cdot\frac{4}{9}}{\frac{1}{2}\cdot\frac{4}{9}
+ \frac{1}{2}\cdot\frac{2}{5}} = \frac{10}{19} \approx 52.6\%$$
\end{subbox}

\begin{subbox}{Wetterprognose (a-priori $\to$ a-posteriori)}
$R$: Prognose schlecht, $S$: schön. $P(R \mid S) = \frac{1}{4}$,
$P(R \mid \overline{S}) = \frac{4}{5}$, Prior $P(S) = \frac{2}{3}$:
$$P(S \mid R) = \frac{\frac{2}{3}\cdot\frac{1}{4}}{\frac{2}{3}\cdot\frac{1}{4}
+ \frac{1}{3}\cdot\frac{4}{5}} = \frac{5}{13} \approx 38.5\%$$
\end{subbox}

\begin{subbox}{Spam-Filter (Bayes-Klassifikator)}
\textbf{1 Wort:} \emph{Rolex} in 250/2000 Spam, 5/1000 Ham,
$P(S) = P(\overline{S}) = \frac{1}{2}$:
$$P(S \mid R) = \frac{0.125}{0.125 + 0.005} = \frac{25}{26} \approx 96.2\%$$
$> 0.9 \Rightarrow$ ablehnen.\\
\textbf{2 Wörter} (Annahme: unabhängig gegeben $S$ bzw.\ $\overline{S}$):
$P(A \cap U \mid S) = P(A \mid S) \cdot P(U \mid S)$.
\emph{Aktie} 400/2000 | 60/1000, \emph{unterbewertet} 200/2000 | 25/1000:
\begin{align*}
P(S \mid A \cap U) &= \frac{0.2 \cdot 0.1}{0.2 \cdot 0.1 + 0.06 \cdot 0.025}\\
&= \frac{40}{43} \approx 93\%
\end{align*}
(Prior $\frac{1}{2}$ kürzt sich weg.)
\end{subbox}

\begin{subbox}{Medizinischer Test (Basisrate!)}
TBC-Rate $0.1\%$, Sensitivität $P(D \mid T) = 0.9$,
Fehlalarm $P(D \mid \overline{T}) = 0.01$:
\begin{align*}
P(T \mid D) &= \frac{0.001 \cdot 0.9}{0.001 \cdot 0.9 + 0.999 \cdot 0.01}\\
&\approx 8.26\%
\end{align*}
Mit Rate $0.19\%$, $P(D \mid T) = 0.98$, Fehlalarm $< 13\cdot10^{-6}$:
$P(T \mid D) > 99.3\%$ (Nenner mit oberer Schranke $\Rightarrow$ untere
Schranke).
\end{subbox}

\begin{subbox}{Verrauschter Kanal mit Paritätsbit}
0110 gesendet, jedes Bit kippt unabhängig mit $p$, $q = 1-p$.
$G$: gerade Anzahl Einsen empfangen ($\Leftrightarrow$ 0, 2 oder 4 Bits
gekippt):
$$P(\text{korrekt} \mid G) = \frac{q^4}{q^4 + \binom{4}{2}q^2p^2 + p^4}$$
$p = 1\%$: $\approx 99.94\%$; \ $p = 10\%$: $\approx 93.09\%$
\end{subbox}

\begin{subbox}{Gestörte Signale}
$0 : 1$ im Verhältnis $5:3$, gestört: $\frac{2}{5}$ der Nullen,
$\frac{1}{3}$ der Einsen. $S_i$: gesendet, $E_i$: empfangen:
$$P(S_0 \mid E_0) = \frac{\frac{5}{8}\cdot\frac{3}{5}}
{\frac{5}{8}\cdot\frac{3}{5} + \frac{3}{8}\cdot\frac{1}{3}} = 75\%$$
$$P(S_1 \mid E_1) = \frac{\frac{3}{8}\cdot\frac{2}{3}}
{\frac{3}{8}\cdot\frac{2}{3} + \frac{5}{8}\cdot\frac{2}{5}} = 50\%$$
\end{subbox}

\begin{subbox}{Faire vs.\ $p$-gefälschte Münze (uniform gezogen)}
$P(G \mid W) = \frac{P(W \mid G)}{P(W \mid G) + P(W \mid F)}$
(Prior $\frac{1}{2}$ kürzt sich):
\begin{itemize}
    \item $p = \frac{2}{3}$, K: $\frac{2/3}{2/3 + 1/2} = \frac{4}{7}
          \approx 57.1\%$
    \item $p = \frac{3}{4}$, KKKZ: $\frac{27/256}{27/256 + 16/256} =
          \frac{27}{43} \approx 62.8\%$
    \item $p = \frac{3}{4}$, KZZZKZ: $\frac{9}{9 + 64} = \frac{9}{73}
          \approx 12.3\%$
\end{itemize}
Nur Anzahl K/Z zählt, nicht die Reihenfolge.
\end{subbox}

\begin{subbox}{Urne raten (2 Kugeln ohne Zurücklegen)}
I: 2 grün, 1 rot; \ II: 4 grün, 1 rot; Urne uniform.\\
I: $P(GG) = P(GR) = P(RG) = \frac{1}{3}$, $P(RR) = 0$\\
II: $P(GG) = \frac{3}{5}$, $P(GR) = P(RG) = \frac{1}{5}$, $P(RR) = 0$\\
\textbf{Plan} (grössere Likelihood raten): GG$\to$II, GR/RG$\to$I:
$$P(\text{richtig}) = \tfrac{1}{2}\left(\tfrac{1}{3} + \tfrac{1}{3}\right)
+ \tfrac{1}{2} \cdot \tfrac{3}{5} = \tfrac{19}{30} \approx 63.3\%$$
$$P(\text{II} \mid GG) = \frac{\frac{1}{2}\cdot\frac{3}{5}}
{\frac{1}{2}\cdot\frac{3}{5} + \frac{1}{2}\cdot\frac{1}{3}} = \frac{9}{14}
\approx 64.3\%$$
\end{subbox}

\begin{subbox}{Häufige Fehler \& Tipps}
\begin{itemize}
    \item $P(A \mid B)$ mit $P(B \mid A)$ verwechseln -- genau das
          korrigiert Bayes.
    \item \textbf{Basisrate} ignorieren: guter Test + seltene Krankheit
          $\Rightarrow$ $P(\text{krank} \mid +)$ trotzdem klein.
    \item Nenner $P(B)$ über \emph{alle} Partitionsteile summieren, auch
          die mit kleinem $P(B \mid \Omega_i)$.
    \item ``x\,\% bleiben unerkannt'' $= P(\overline{D} \mid T)$
          $\Rightarrow$ $P(D \mid T) = 1 - x$.
    \item Uniformer Prior kürzt sich: dann nur Likelihoods vergleichen.
\end{itemize}
\end{subbox}

% ============================================================================
% TIPS FOR AVOIDING OVERFULL BOXES
% ============================================================================
%
% 1. Keep formulas concise - split long equations over multiple lines
% 2. Use shorter variable names in inline math
% 3. Abbreviate table headers
% 4. Use align* environment for multi-line equations
% 5. Test compile frequently to catch overflow early
%
% Example of splitting a long formula:
% Instead of: $$a + b + c + d + e = f + g + h + i + j$$
% Use:
% \begin{align*}
% a + b + c + d + e &= f + g + h \\
%                   &\quad + i + j
% \end{align*}
%
% ============================================================================

\end{document}